% Generated draft: requires mathematical and source review.
\documentclass[11pt]{article}
\usepackage[margin=25mm]{geometry}
\usepackage{fontspec}
\setmainfont{Latin Modern Roman}
\usepackage{amsmath,amssymb,mathtools,unicode-math}
\setmathfont{Latin Modern Math}
\usepackage{enumitem,needspace,xurl,hyperref}
\setlength{\parindent}{0pt}
\setlength{\parskip}{5pt}
\setlength{\emergencystretch}{3em}
\newcommand{\N}{\mathbb{N}}
\newcommand{\Z}{\mathbb{Z}}
\newcommand{\Q}{\mathbb{Q}}
\newcommand{\R}{\mathbb{R}}
\newcommand{\C}{\mathbb{C}}
\newcommand{\F}{\mathbb{F}}
\newcommand{\A}{\mathbb{A}}
\newcommand{\PP}{\mathbb{P}}
\newcommand{\E}{\mathbb{E}}
\newcommand{\eps}{\varepsilon}
\newcommand{\dd}{\,\mathrm{d}}
\begin{document}
\section*{UnsolvedMath ID: 701553}
\textbf{Source problem code:} catalogue entry
\subsection{Definitions and mathematical statement}
Let \(I\) be an infinite index set, and let \((A_i)_{i\in I}\) be a family of finitely generated abelian groups. Let \(H\) be a group equipped with injective homomorphisms \(\iota_i:H\hookrightarrow A_i\) for every \(i\in I\). Write
\[
G=\mathop{\ast}_{H,\,i\in I} A_i
\]
for their free product with amalgamation, identifying the copies \(\iota_i(H)\) of \(H\). Identify \(H\) with its canonical image in \(G\). Since \(H\) embeds in an abelian group, it is abelian, and its torsion subgroup is
\[
\operatorname{Tor}(H)=\{h\in H:\text{there exists an integer }n\geq 1\text{ with }h^n=1\}.
\]
The source asks whether, for every such amalgamated free product,
\[
\psi(G)=\operatorname{Tor}(H).
\]
However, the supplied record does not define \(\psi(G)\). Thus this equality records the intended question only schematically; its precise mathematical meaning requires the source's definition of \(\psi\). No countability restriction on \(I\), or properness condition on the embeddings of \(H\), is stated in the supplied record.
\subsection{Short English statement}
M. K. Azarian asks whether, for a free product of infinitely many finitely generated abelian groups amalgamating a common subgroup \(H\), the object denoted by \(\psi(G)\) equals the torsion subgroup of \(H\). The supplied record leaves \(\psi\) undefined, so the question requires clarification before it is a fully specified mathematical statement.
\subsection{Sources}
\begin{itemize}
\item Local UnsolvedMath registry, permanent ID \(701553\), exact source problem code \texttt{catalogue entry}, title ``Kourovka Notebook Problem 19.5.'' The supplied registry labels the problem as open; this label is not independently verified here.
\item Supplied attribution: M. K. Azarian, proposed in 2018; \emph{The Kourovka Notebook (Unsolved Problems in Group Theory)}, original problem number 19.5. The supplied background names the 19th issue (2018) and the 21st issue edition (2026), with identifier \texttt{arXiv:1401.0300v48}. The collection description names E. I. Khukhro and V. D. Mazurov as editors and describes coverage of Issues 1--21 (1965--2026). These bibliographic details are supplied attribution, not independently checked source contents.
\item Supplied registry update date: \texttt{2026-10-16T14:00:00Z}; upstream revision: \texttt{d371ee917ac9647fbb06cd6fc529deec6f3dda82}. No source URL, literature-source links, or literature-verification date is supplied.
\end{itemize}
\end{document}
